\documentclass[10pt,a4paper]{amsart}
\usepackage{amsmath,amssymb,amsthm,graphicx}
\usepackage{verbatim}
\newenvironment{reference}{\comment}{\endcomment}
%\newenvironment{reference}{}{}
\newenvironment{slogan}{\comment}{\endcomment}
\newenvironment{history}{\comment}{\endcomment}

% For commutative diagrams we use Xy-pic
\usepackage[all]{xy}

% We use 2cell for 2-commutative diagrams.
\xyoption{2cell}
% Theorem environments.
%
\theoremstyle{plain}
\newtheorem{theorem}[subsection]{Theorem}
\newtheorem{proposition}[subsection]{Proposition}
\newtheorem{lemma}[subsection]{Lemma}

\theoremstyle{definition}
\newtheorem{definition}[subsection]{Definition}
\newtheorem{example}[subsection]{Example}
\newtheorem{exercise}[subsection]{Exercise}
\newtheorem{situation}[subsection]{Situation}

\theoremstyle{remark}
\newtheorem{remark}[subsection]{Remark}
\newtheorem{remarks}[subsection]{Remarks}

\numberwithin{equation}{subsection}

% Macros
%
\def\lim{\mathop{\mathrm{lim}}\nolimits}
\def\colim{\mathop{\mathrm{colim}}\nolimits}
\def\Spec{\mathop{\mathrm{Spec}}}
\def\Hom{\mathop{\mathrm{Hom}}\nolimits}
\def\Ext{\mathop{\mathrm{Ext}}\nolimits}
\def\SheafHom{\mathop{\mathcal{H}\!\mathit{om}}\nolimits}
\def\SheafExt{\mathop{\mathcal{E}\!\mathit{xt}}\nolimits}
\def\Sch{\mathit{Sch}}
\def\Mor{\mathop{\mathrm{Mor}}\nolimits}
\def\Ob{\mathop{\mathrm{Ob}}\nolimits}
\def\Sh{\mathop{\mathit{Sh}}\nolimits}
\def\NL{\mathop{N\!L}\nolimits}
\def\CH{\mathop{\mathrm{CH}}\nolimits}
\def\proetale{{pro\text{-}\acute{e}tale}}
\def\etale{{\acute{e}tale}}
\def\QCoh{\mathit{QCoh}}
\def\Ker{\mathop{\mathrm{Ker}}}
\def\Im{\mathop{\mathrm{Im}}}
\def\Coker{\mathop{\mathrm{Coker}}}
\def\Coim{\mathop{\mathrm{Coim}}}

% Boxtimes
%
\DeclareMathSymbol{\boxtimes}{\mathbin}{AMSa}{"02}

%
% Macros for moduli stacks/spaces
%
\def\QCohstack{\mathcal{QC}\!\mathit{oh}}
\def\Cohstack{\mathcal{C}\!\mathit{oh}}
\def\Spacesstack{\mathcal{S}\!\mathit{paces}}
\def\Quotfunctor{\mathrm{Quot}}
\def\Hilbfunctor{\mathrm{Hilb}}
\def\Curvesstack{\mathcal{C}\!\mathit{urves}}
\def\Polarizedstack{\mathcal{P}\!\mathit{olarized}}
\def\Complexesstack{\mathcal{C}\!\mathit{omplexes}}
% \Pic is the operator that assigns to X its picard group, usage \Pic(X)
% \Picardstack_{X/B} denotes the Picard stack of X over B
% \Picardfunctor_{X/B} denotes the Picard functor of X over B
\def\Pic{\mathop{\mathrm{Pic}}\nolimits}
\def\Picardstack{\mathcal{P}\!\mathit{ic}}
\def\Picardfunctor{\mathrm{Pic}}
\def\Deformationcategory{\mathcal{D}\!\mathit{ef}}
\usepackage{cleveref}
\crefname{originallemma}{original lemma}{original lemmas}
\Crefname{originallemma}{Original lemma}{Original lemmas}
\makeatletter
\newlabel{dualizing-lemma-finite}{{08YX}{source}{}{sourcepointer-1}{}}
\newlabel{dualizing-lemma-finite@cref}{{[originallemma][0][]08YX}{[1][0][]source}{}{}{}}
\makeatother
\crefname{originallemma}{original lemma}{original lemmas}
\Crefname{originallemma}{Original lemma}{Original lemmas}
\makeatletter
\newlabel{dualizing-lemma-union-artinian}{{08Z4}{source}{}{sourcepointer-2}{}}
\newlabel{dualizing-lemma-union-artinian@cref}{{[originallemma][0][]08Z4}{[1][0][]source}{}{}{}}
\makeatother
\crefname{originalproposition}{original proposition}{original propositions}
\Crefname{originalproposition}{Original proposition}{Original propositions}
\makeatletter
\newlabel{dualizing-proposition-structure-injectives-noetherian}{{08YA}{source}{}{sourcepointer-3}{}}
\newlabel{dualizing-proposition-structure-injectives-noetherian@cref}{{[originalproposition][0][]08YA}{[1][0][]source}{}{}{}}
\makeatother
\title{Injective quasi-coherent modules on an affine Noetherian scheme are filtered colimits of Artinian pushforwards}
\author{The Stacks Project Authors}
\date{stacks-project:spaces-perfect.tex:local-snapshot-20261012; GNU FDL 1.2 or later}
\begin{document}
\maketitle
\noindent\textit{Editorial note.} All delivered mathematical text is the author's own, copied from the pinned Stacks Project chapter file spaces-perfect.tex; no statement or proof step was rewritten, reordered or removed.\par\medskip
\noindent\textit{Editorial note.} The counted proof is the author's complete argument: the proof reduces to indecomposable injectives via the structure theorem, writes the injective hull as the union of its Artinian layers using the union lemma, and reads the layers as pushforwards of coherent sheaves from Artinian spectra\par\medskip
\noindent\textit{Editorial note.} The label lemma-affine-injective-colimit-direct-sum-pushforwards-artin names the counted statement, and the proof environment immediately following it in the source is delivered verbatim as its proof.\par\medskip
\noindent\textit{Editorial note.} The counted statement is carried by the author's own lemma environment in the chapter source.\par\medskip
\noindent\textit{Allowed standard prerequisites (editorial).}\begin{itemize}
\item the equivalence between QCoh(O\_S) on an affine scheme and the category of A-modules
\item the structure theorem for injective modules over a Noetherian ring: every injective is a direct sum of indecomposable injectives, each the injective hull of a residue field (Dualizing Complexes, Proposition dualizing-proposition-structure-injectives-noetherian)
\item the description of the injective hull of kappa(p) as the union of its Artinian submodules J[p\textasciicircum{}n], the injective hulls over A\_p/p\textasciicircum{}n A\_p (Dualizing Complexes, Lemma dualizing-lemma-union-artinian)
\item finiteness of the pieces (Dualizing Complexes, Lemma dualizing-lemma-finite)
\end{itemize}\medskip
\ifdefined\hypertarget\hypertarget{sourcepointer-1}{}\fi
\noindent\textit{Source reference.} Original lemma 08YX: sources/stacks-chow/tags.txt, source lines 11355--11355. Prerequisite pointer invoked at proof line(s) 2864; it names a standard fact of a sibling chapter that the counted argument uses, and every step of the counted proof built on it is retained verbatim in the delivered body.\par\medskip
\ifdefined\hypertarget\hypertarget{sourcepointer-2}{}\fi
\noindent\textit{Source reference.} Original lemma 08Z4: sources/stacks-chow/tags.txt, source lines 11362--11362. Prerequisite pointer invoked at proof line(s) 2857; it names a standard fact of a sibling chapter that the counted argument uses, and every step of the counted proof built on it is retained verbatim in the delivered body.\par\medskip
\ifdefined\hypertarget\hypertarget{sourcepointer-3}{}\fi
\noindent\textit{Source reference.} Original proposition 08YA: sources/stacks-chow/tags.txt, source lines 11333--11333. Prerequisite pointer invoked at proof line(s) 2848; it names a standard fact of a sibling chapter that the counted argument uses, and every step of the counted proof built on it is retained verbatim in the delivered body.\par\medskip
\begin{lemma}
\label{lemma-affine-injective-colimit-direct-sum-pushforwards-artin}
Let $S$ be a Noetherian affine scheme. Every injective object of
$\QCoh(\mathcal{O}_S)$ is a filtered colimit $\colim_i \mathcal{F}_i$
of quasi-coherent sheaves of the form
$$
\mathcal{F}_i = (Z_i \to S)_*\mathcal{G}_i
$$
where $Z_i$ is the spectrum of an Artinian ring and $\mathcal{G}_i$
is a coherent module on $Z_i$.
\end{lemma}

\begin{proof}
Let $S = \Spec(A)$. Let $\mathcal{J}$ be an injective object of
$\QCoh(\mathcal{O}_S)$. Since $\QCoh(\mathcal{O}_S)$ is
equivalent to the category of $A$-modules we see that $\mathcal{J}$
is equal to $\widetilde{J}$ for some injective $A$-module $J$.
By Dualizing Complexes, Proposition
\ref{dualizing-proposition-structure-injectives-noetherian}
we can write $J = \bigoplus E_\alpha$ with $E_\alpha$ indecomposable
and therefore isomorphic to the injective hull of a reside field
at a point. Thus (because finite disjoint unions of Artinian schemes
are Artinian) we may assume that $J$ is the injective hull
of $\kappa(\mathfrak p)$ for some prime $\mathfrak p$ of $A$.
Then $J = \bigcup J[\mathfrak p^n]$ where $J[\mathfrak p^n]$ is
the injective hull of $\kappa(\mathfrak p)$ over
$A_\mathfrak/\mathfrak p^nA_\mathfrak p$, see
Dualizing Complexes, Lemma \ref{dualizing-lemma-union-artinian}.
Thus $\widetilde{J}$ is the colimit of the sheaves
$(Z_n \to X)_*\mathcal{G}_n$ where
$Z_n = \Spec(A_\mathfrak p/\mathfrak p^nA_\mathfrak p)$ and
$\mathfrak G_n$ the coherent sheaf associated to the
finite $A_\mathfrak/\mathfrak p^nA_\mathfrak p$-module $J[\mathfrak p^n]$.
Finiteness follows from
Dualizing Complexes, Lemma \ref{dualizing-lemma-finite}.
\end{proof}

\end{document}
